\section{Related Work}
\label{sec:related}

\subsection{Coastal water-level forecasting}
Traditional coastal water-level forecasting combines deterministic tidal
constituents with empirical or dynamical storm-surge models~\citep{yang2023storm}.
In recent years, neural approaches have improved forecast quality on held-out
gauges. \citet{sergeev2024spatial} demonstrated that information from multiple
gauges along a coastline substantially improves forecast accuracy for downstream
stations. \citet{tong2024dual} introduced dual-timescale residual memory that
separately models slow tidal evolution and fast storm-driven residuals, with
particular attention to extreme events. These works establish that (a) the
de-tided residual signal carries the most operationally relevant information during
storms, and (b) extreme events demand model architectures that differ from those
optimized for benign conditions.

\subsection{Data augmentation for time-series forecasting}
Data augmentation for time series is reviewed by~\citet{hoppe2024diffusion} and
related surveys. Standard approaches include: (1) \emph{oversampling} rare events
or weighting them in the loss function, (2) \emph{perturbation} with additive or
multiplicative noise, (3) \emph{interpolation} between adjacent real windows, and
(4) \emph{generative models} including GANs, VAEs, and diffusion. A key
distinction for coastal water levels is whether augmentation preserves the
\emph{physical plausibility} of the residual signal: unrealistic samples may
actually degrade downstream forecasting by corrupting the training distribution.

\subsection{Conditional diffusion models for time series}
Conditional score-based diffusion models have been applied to time-series
imputation~\citep{tashiro2022csdi,ma2023marginal} and generation. The canonical
formulation~\citep{ho2020denoising} defines a forward process
$q(x_t | x_0) = \mathcal{N}(\sqrt{\bar{\alpha}_t} x_0, (1 - \bar{\alpha}_t) I)$
and trains a network $\epsilon_\theta(x_t, t, c)$ to predict the injected noise,
where $c$ denotes conditioning information. During generation, the reverse process
iteratively denoises from $T$ back to 0:
\begin{equation}
    x_{t-1} = \frac{1}{\sqrt{\alpha_t}} \left( x_t - \frac{\beta_t}{\sqrt{1-\bar{\alpha}_t}} \epsilon_\theta(x_t, t, c) \right) + \sigma_t z,
    \label{eq:reverse}
\end{equation}
where $\sigma_t^2 = \beta_t (1 - \bar{\alpha}_{t-1}) / (1 - \bar{\alpha}_t)$ is the
posterior variance and $z \sim \mathcal{N}(0, I)$.

\subsection{This work's niche}
To our knowledge, this is the first systematic application of conditional diffusion
augmentation to coastal water-level forecasting with (a) event-grouped leakage-safe
splits, (b) train-only threshold estimation, (c) matched-budget comparison across
augmentation methods, and (d) pre-registration of hypotheses before running
experiments.
